\documentclass[11pt]{article}

\usepackage[margin=1in]{geometry}
\usepackage{amsmath, amssymb}
\usepackage{lmodern}

\setlength{\parindent}{0pt}
\setlength{\parskip}{0.8em}

\begin{document}

\begin{center}
{\LARGE Raw p-value Method}
\end{center}

\section*{Definition}

The raw p-value method is a mass-univariate baseline. For each time point and
channel, we perform a separate one-sample t-test across participants on the
condition difference.

For each time point \(t\) and channel \(c\), the null hypothesis is:

\[
H_0: \Delta EEG_{\mathrm{participant}}(t, c) = 0.
\]

The p-value at each time-channel cell is:

\[
p_{t,c}
=
p\text{-value from a t-test across participants at time }t
\text{ and channel }c.
\]

The significant cells are defined as:

\[
\mathrm{raw\_p\ mask}_{t,c}
=
\mathbf{1}(p_{t,c} < \alpha).
\]

\end{document}
